% Exemplo: elementos básicos das fórmulas (^, _, chaves, \frac, \sqrt, \sum, \int, letras gregas)
% Compilar: latexmk -pdf elementos.tex             Testado com: TeX Live 2024 ou mais recente
% Cada trecho fica em uma página. É o ponto de partida antes de modos.tex (destaque e numeração).
\documentclass[a4paper]{article}
\usepackage[T1]{fontenc}
\usepackage{lmodern}
\usepackage[brazilian]{babel}
\usepackage{amsmath}

\begin{document}

%%BEGIN:elementosbasicos
Expoente com \verb|^| e índice com \verb|_|: $x^2$, $x_i$ e $x_i^2$. Mais
de um caractere vai entre chaves: $e^{-x}$ e $a_{ij}$, mas $e^-x$ e
$a_ij$ mostram o que acontece sem elas.

Frações e raízes: $\frac{a}{b}$, $\frac{1}{1+x^2}$, $\sqrt{2}$ e
$\sqrt[3]{x}$. Em destaque, a fração fica em tamanho normal:
\[
  \bar{x} = \frac{1}{n} \sum_{i=1}^{n} x_i
  \qquad
  \int_{0}^{1} x^2 \, dx = \frac{1}{3}
\]
No somatório e na integral, \verb|_| e \verb|^| dão os limites; no
texto eles vão ao lado: $\sum_{i=1}^{n} x_i$.
%%END:elementosbasicos

\clearpage

%%BEGIN:letrasgregas
Letras gregas são comandos com o nome da letra: $\alpha$, $\beta$,
$\gamma$, $\delta$, $\varepsilon$, $\lambda$, $\mu$, $\pi$, $\sigma$ e
$\omega$. A inicial maiúscula dá a letra maiúscula, quando ela difere da
latina: $\Gamma$, $\Delta$, $\Sigma$ e $\Omega$. Assim, a média e o
desvio-padrão de uma população são $\mu$ e $\sigma$, e a variância é
$\sigma^2$.
%%END:letrasgregas

\end{document}
